% Abstract text (150--250 words). Keywords are set in main.tex.
Las clasificaciones de autómatas celulares agrupan reglas por su comportamiento, pero no miden qué tan similares son. Buscamos un índice de similitud basado en la evolución que permita comparar reglas, incluso con distinto número de estados o de vecinos. Primero definimos cuatro relaciones exactas (permutación de estados, espejo, reducción de estados y reducción de vecindad), que reducen las 256 reglas de los autómatas celulares elementales (ECA) a 81 clases. Después comparamos el grafo de transición en un anillo de $n$ células: su forma canónica, junto con el cociente por traslaciones, distingue las 88 clases por permutación y espejo con solo 6 células, pero la comparación es binaria. Como medida gradual proponemos los momentos espectrales con masa de atractor, un vector de longitud fija que, a diferencia de los eigenvalores y los momentos espectrales, cuenta cuántas configuraciones llegan a cada atractor y en cuántos pasos. Con la distancia coseno, este vector distingue 86 de las 88 clases (87 junto con el cociente), incluidas las reglas 0 y 8, y da distancia cero a las relaciones que conservan el grafo y, en el cociente, a las que difieren por una traslación, pero no identifica la reducción de estados. Además, depende del tamaño del espacio y agrupa las ECA en cuatro familias según el periodo dominante de sus atractores.
